\section{Introduction}\label{sec:introduction}
Let $N\ge2$ be an integer.
\par
In \cite{Kashaev:MODPLA95}, R.~Kashaev introduced a link invariant $\langle K\rangle_{N}$ for a knot $K$ in the three-sphere $S^3$ by using the so-called quantum dilogarithm.
In \cite{Kashaev:LETMP97}, he studied its asymptotic behavior when $N\to\infty$ for several hyperbolic knots and conjectured that it would determine its hyperbolic volume for any {\em hyperbolic} knot, where a knot is called hyperbolic if its complement possesses a complete hyperbolic structure with finite volume.
More precisely he conjectured
\begin{conj}[Kashaev's conjecture]
Let $H$ be a {\em hyperbolic} knot in $S^3$.
Then the following equality holds.
\begin{equation}\label{eq:Kashaev}
  \lim_{N\to\infty}
  \frac{\log\left|\langle H\rangle_{N}\right|}{N}
  =
  \frac{\operatorname{V}\left(S^3\setminus{H}\right)}{2\pi},
\end{equation}
where $\operatorname{V}\left(S^3\setminus{H}\right)$ is the hyperbolic volume.
\end{conj}
\par
Let $J_{N}(K;q)$ be the colored Jones polynomial associated with the $N$-dimensional irreducible representation of the Lie algebra $\sl(2;\C)$ \cite{Kirby/Melvin:INVEM1991,Reshetikhin/Turaev:INVEM1991}.
We normalize it so that $J_{N}(U;q)=1$ for the unknot $U$, and that $J_{2}(K;q)$ is the ordinary Jones polynomial \cite{Jones:BULAM31985}.
In \cite{Murakami/Murakami:ACTAM12001}, J.~Murakami and the first author proved that Kashaev's invariant $\langle K\rangle_{N}$ coincides with the $N$-dimensional colored Jones polynomial $J_{N}(K;q)$  evaluated at $q=e^{2\pi\i/N}$.
They also generalized Kashaev's conjecture for general knots:
\begin{conj}[Volume Conjecture]
For {\em any} knot $K\subset{S^3}$, we have
\begin{equation}\label{eq:VC}
  \lim_{N\to\infty}
  \frac{\log\left|J_{N}\left(K;e^{2\pi\i/N}\right)\right|}{N}
  =
  \frac{\Vol\left(S^3\setminus{K}\right)}{2\pi}.
\end{equation}
Here, $\Vol(S^3\setminus{K})$ is the simplicial volume, also known as the Gromov norm \cite{Gromov:INSHE82} {\rm(}see also \cite{Soma:INVEM1981}{\rm)}, which is normalized so that $\Vol\left(S^3\setminus{H}\right)=\operatorname{V}\left(S^3\setminus{H}\right)$ if $H$ is hyperbolic.
\end{conj}
\par
Kashaev's conjecture was complexified by T.~Takata, J.~Murakami, M.~Okamoto, Y.~Yokota and the first author in \cite{Murakami/Murakami/Okamoto/Takata/Yokota:EXPMA02} by dropping the absolute value symbol (and replacing $\langle H\rangle_{N}$ with $J_{N}\left(H;e^{2\pi\i/N}\right)$) in \eqref{eq:Kashaev}.
\begin{conj}[Complexification of Kashaev's conjecture]
For a hyperbolic knot $H$, we have
\begin{equation}\label{eq:CVC}
  \lim_{N\to\infty}
  \frac{\log J_{N}\left(H;e^{2\pi\i/N}\right)}{N}
  =
  \frac{\operatorname{V}\left(S^3\setminus{H}\right)+\i\CS(S^3\setminus{H})}{2\pi},
\end{equation}
where $\CS$ is the $SO(3)$ Chern--Simons invariant of $H$ associated with the Levi-Civita connection \cite{Meyerhoff:LMSLN112}.
\end{conj}
\par
More detailed asymptotic formulas are known for several knots \cite{Andersen/Hansen:JKNOT2006,Ohtsuki:QT2016,Ohtsuki/Yokota:MATPC2018}:
\begin{equation*}
  \langle H\rangle_{N}
  \underset{N\to\infty}{\sim}
  \omega(H)N^{3/2}\exp\left(\frac{N}{2\pi}\operatorname{CV}(H)\right),
\end{equation*}
where $\operatorname{CV}(H):=\Vol\left(S^3\setminus{H}\right)+\i\CS(S^3\setminus{H})$ is the complex volume and $2\i\omega(H)^2$ is the adjoint (cohomological) Reidemeister torsion twisted by the holonomy representation of $\pi_1(S^3\setminus{H})$ to $\SL(2;\C)$ \cite{Ohtsuki/Takata:GEOTO2015}.
\par
It was also generalized by replacing $2\pi\i$ with another complex number $\xi$ \cite{Dimofte/Gukov/Lenells/Zagier:CNTP2010,Murakami:JTOP2013}.
\begin{conj}[Generalized Volume Conjecture]
Let $\xi\ne2\pi\i$ be a complex number close to $2\pi\i$, and $H$ a hyperbolic knot.
Then we have
\begin{equation}\label{eq:GVC}
  J_{N}\left(H;e^{\xi/N}\right)
  \underset{N\to\infty}{\sim}
  \frac{C}{\sinh(\xi)}\times\tau(\xi)^{1/2}\left(\frac{N}{\xi}\right)^d
  \exp\left(\frac{N}{\xi}S(\xi)\right)
\end{equation}
with $d$ a constant.
Here $S(\xi)$ and $\tau(\xi)$ are related to the $\SL(2;\C)$ Chern--Simons invariant and the adjoint Reidemeister torsion, respectively, both associated with a certain representation of $\pi_1(S^3\setminus{H})$ to $\SL(2;\C)$.
\end{conj}
\par
In this paper, we are mainly interested in the figure-eight knot.
We first list known results for the asymptotic behavior of the colored Jones polynomial of the figure-eight knot $E$.
\par
We define a function $\varphi$ as
\begin{equation*}
  \varphi(\xi)
  :=
  \arccosh\left(\cosh(\xi)-\frac{1}{2}\right)
\end{equation*}
for a complex number $\xi$, where we use the following branch of $\arccosh$:
\begin{equation*}
  \arccosh(x)
  :=
  \log\left(x-\i\sqrt{1-x^2}\right),
\end{equation*}
and we choose the branch cut of $\log$ as $(-\infty,0)$.
Note that $\varphi(0)=-\pi\i/3$ and $\varphi(\kappa)=0$, where $\kappa:=\log\left(\frac{3+\sqrt{5}}{2}\right)$.
We also define
\begin{align*}
  S(\xi)
  &:=
  \Li_2(e^{-\xi-\varphi(\xi)})-\Li_2(e^{-\xi+\varphi(\xi)})+\xi\varphi(\xi),
  \\
  \tilde{S}(\xi)
  &:=
  \Li_2(e^{-\xi-\varphi(\xi)})-\Li_2(e^{-\xi+\varphi(\xi)})
  +(\xi-2\pi\i)(\varphi(\xi)+2\pi\i),
  \\
  T(\xi)
  &:=
  \frac{2}{\sqrt{(2\cosh(\xi)+1)(2\cosh(\xi)-3)}},
\end{align*}
where
\begin{equation*}
  \Li_2(z):=-\int_{0}^{z}\frac{\log(1-x)}{x}\,dx
\end{equation*}
is the dilogarithm function, where we choose the branch cut as $(1,\infty)$.
Note that $S(0)=\tilde{S}(2\pi\i)=\Li_2(e^{\pi\i/3})-\Li_2(e^{-\pi\i/3})=\i\times 2.02988...$.
\par
The quantities $S(\xi)$, $\tilde{S}(\xi)$, and $T(\xi)$ are related to the Chern--Simons invariant and the adjoint Reidemeister torsion associated with a certain representation of $\pi_1\left(S^3\setminus{E}\right)$ to $\SL(2;\C)$.
See Section~\ref{sec:CST} for details.
\begin{enumerate}
\item
$\xi=2\pi\i$:
\par
This corresponds to the case of the Volume Conjecture.
Kashaev sketched a proof of \eqref{eq:VC} in \cite{Kashaev:LETMP97}, and then T.~Ekholm gave a detailed proof of \eqref{eq:VC}.
See for example \cite[Section~3.2]{Murakami/Yokota:2018} for Ekholm's proof.
Later, J.E.~Andersen and S.K.~Hansen \cite[Theorem~1]{Andersen/Hansen:JKNOT2006} followed Kashaev's method to prove the following asymptotic formula.
\begin{equation*}
  J_{N}\left(E;e^{2\pi\i/N}\right)
  \underset{N\to\infty}{\sim}
  2\pi^{3/2}\left(\frac{N}{2\pi\i}\right)^{3/2}T(2\pi\i)^{1/2}
  \exp\left(\frac{N}{2\pi\i}\tilde{S}(2\pi\i)\right).
\end{equation*}
Since $S(0)=\i\Vol\left(S^3\setminus{E}\right)$, this refines the Volume Conjecture for the figure-eight knot.
\item
$\xi$ is close to $2\pi\i$ and not purely imaginary:
\par
Y.~Yokota and the first author proved the following formula \cite{Murakami/Yokota:JREIA2007}.
\begin{equation}\label{eq:Yokota}
  \lim_{N\to\infty}
  \frac{\log J_{N}(E;e^{\xi/N})}{N}
  =
  \frac{\tilde{S}(\xi)}{\xi}.
\end{equation}
\item
$\xi$ is purely imaginary with $5\pi/3<|\xi|<7\pi/3$, and $2\pi/|\xi|$ is irrational with finite irrationality measure:
\par
In \cite{Murakami:KYUMJ2004} (see \cite{Murakami:KYUMJ2016} for correction), the first author proved the following formula.
\begin{equation*}
  \lim_{N\to\infty}
  \frac{\log J_{N}(E;e^{\xi/N})}{N}
  =
  \frac{\tilde{S}(\xi)}{\xi}.
\end{equation*}
See also \cite[6.2.2]{Murakami:Novosibirsk}.
\item
$\xi$ is of the form $2\pi\i+u$ for a real number $u$ with $0<|u|<\kappa$:
\par
The first author proved the following asymptotic formula \cite[Theorem~1.4]{Murakami:JTOP2013}:
\begin{equation}\label{eq:J_top}
  J_{N}\left(E;e^{\xi/N}\right)
  \underset{N\to\infty}{\sim}
  \frac{\sqrt{-\pi}}{2\sinh(u/2)}T(\xi)^{1/2}\left(\frac{N}{\xi}\right)^{1/2}
  \exp\left(\frac{N}{\xi}\tilde{S}(\xi)\right).
\end{equation}
Note that this refines \eqref{eq:Yokota} when $\xi$ is as above.
\item
$\xi$ satisfies the inequalities $|2\cosh{\xi}-2|<1$ and $|\Im{\xi}|<\pi/3$:
\par
In this case, $J_{N}\left(E;e^{\xi/N}\right)$ converges.
In fact, we have
\begin{equation*}
  \lim_{N\to\infty}
  J_{N}\left(E;e^{\xi/N}\right)
  =
  \frac{1}{\Delta(E;e^{\xi})},
\end{equation*}
where $\Delta(E;t)=-t+3-t^{-1}$ is the normalized Alexander polynomial of the figure-eight knot $E$
\cite[Theorem~1.1]{Murakami:JPJGT2007}.
\par
See \cite{Garoufalidis/Le:GEOTO2011} for general knots.
\item
$\xi=\pm\kappa$:
\par
In this case, $J_{N}\left(E;e^{\kappa/N}\right)$ grows polynomially.
More precisely, we have
\begin{equation*}
  J_{N}\left(E;e^{\kappa/N}\right)
  \underset{N\to\infty}{\sim}
  \frac{\Gamma(1/3)}{(3\kappa)^{2/3}}N^{2/3},
\end{equation*}
where $\Gamma(z)$ is the Gamma function \cite[Theorem~1.1]{Hikami/Murakami:COMCM2008}.
\item
$\xi$ is real and $|\xi|>\kappa$:
\par
The first author proved the following formula \cite[Theorem~8.1]{Murakami:KYUMJ2004} (see also \cite[Theorem~3.2]{Murakami:ACTMV2008} and \cite[Lemma~6.7]{Murakami:Novosibirsk}).
\begin{equation}\label{eq:Kyungpook}
  \lim_{N\to\infty}
  \frac{\log J_{N}(E;e^{\xi/N})}{N}
  =
  \frac{S(\xi)}{|\xi|}.
\end{equation}
\end{enumerate}
\par
The purpose of this paper is to refine \eqref{eq:Kyungpook}.
We will show
\begin{thm}\label{thm:main}
If $\xi$ is real and $|\xi|>\kappa$, then we have
\begin{equation*}
  J_{N}\left(E;e^{\xi/N}\right)
  \underset{N\to\infty}{\sim}
  \frac{\sqrt{\pi}}{2\sinh(|\xi|/2)}T(\xi)^{1/2}\left(\frac{N}{|\xi|}\right)^{1/2}
  \exp\left(\frac{N}{|\xi|}S(\xi)\right).
\end{equation*}
\end{thm}
Moreover, we can show that $T(\xi)$ is the Reidemeister torsion twisted by the adjoint action of a representation $\rho$ of $\pi_1(S^3\setminus{E})$ to $\SL(2;\C)$ determined by $\xi$, and that $S(\xi)-\xi\eta/2$ is the Chern--Simons invariant of $\rho$ associated with the meridian and the preferred longitude of $E\subset S^{3}$.
See Section~\ref{sec:CST} for details.
\begin{rem}
Suppose that $u:=\xi-2\pi\i$ is a small complex number or a real number with $|u|<\kappa$.
In this case $\Im\varphi(u)<0$.
By using a well-known formula
\begin{equation*}
  \Li_2(z)+\Li_2(z^{-1})+\frac{\pi^2}{6}+\frac{1}{2}\left(\log(-z)\right)^2=0.
\end{equation*}
we have
\begin{equation*}
\begin{split}
  &\tilde{S}(\xi)
  \\
  =&
  -\Li_2(e^{\xi+\varphi(\xi)})+\Li_2(e^{\xi-\varphi(\xi)})
  -\frac{1}{2}\left(\log(-e^{\xi+\varphi(\xi)})\right)^2
  +\frac{1}{2}\left(\log(-e^{\xi-\varphi(\xi)})\right)^2
  \\
  &+u(\varphi(\xi)+2\pi\i)
  \\
  \\
  =&
  -\Li_2(e^{u+\varphi(u)})+\Li_2(e^{u-\varphi(u)})
  -\frac{1}{2}\left(u+\varphi(u)+\pi\i\right)^2
  \\
  &+\frac{1}{2}\left(u-\varphi(u)-\pi\i\right)^2
   +(2\pi\i+u)\varphi(u)
  \\
  =&
  -\Li_2(e^{u+\varphi(u)})+\Li_2(e^{u-\varphi(u)})-u\varphi(u).
\end{split}
\end{equation*}
In the second equality, we use the fact that $\Im\varphi(u)<0$.
Therefore \eqref{eq:J_top} coincides with the formula appearing in \cite[Theorem~1.4]{Murakami:JTOP2013}.
\end{rem}
\begin{rem}
In \cite{Murakami:JTOP2013}, the first author followed \cite{Andersen/Hansen:JKNOT2006} to obtain the asymptotic formula, but in the current paper we follow \cite{Ohtsuki:QT2016}.
\end{rem}